\documentclass[UTF8]{ctexart}

% 支持从项目根目录、强化学习目录或当前文档目录读取共享样式。
\makeatletter
\def\input@path{{../}{../../}}
\makeatother
\usepackage{researchnotes}

\title{贝尔曼方程中回报与价值的替换}
\author{}
\date{}

\begin{document}
\maketitle

\keyterm{这里并不是说 \(G_{t+1}=V(S_{t+1})\)，而是说：在外面的条件期望中，可以用后者替换前者，平均结果不变。} \(G_{t+1}\) 是从下一时刻开始的随机回报；\(V(s')\) 是从状态 \(s'\) 出发时，这个回报的平均值。例如，从状态 \(A\) 出发，回报可能各以一半的概率取 \(0\) 或 \(10\)，那么 \(V(A)=5\)。实际回报是 \(0\) 或 \(10\)，并不等于 \(5\)，但它的期望等于 \(5\)。

再把“下一步到哪个状态”也考虑进来。假设当前在 \(s\)，下一步各以一半的概率到达 \(A\) 或 \(B\)：到达 \(A\) 后，后续回报等概率为 \(0\) 或 \(10\)，所以 \(V(A)=5\)；到达 \(B\) 后，后续回报等概率为 \(2\) 或 \(4\)，所以 \(V(B)=3\)。此时，\keyterm{直接对所有可能的回报求平均，与先计算各状态内的平均回报、再对状态求平均，结果相同：}
\[
\underbrace{\mathbb E[G_{t+1}\mid S_t=s]}_{\text{直接平均四种可能的回报}}
=\frac14(0+10+2+4)=4,
\qquad
\underbrace{\mathbb E[V(S_{t+1})\mid S_t=s]}_{\text{平均两个状态的价值}}
=\frac12\times5+\frac12\times3=4.
\]

图中省略的正是这个“分组求平均”的过程，数学上叫作\keyterm{全期望公式}。沿用图中的时间齐次马尔可夫奖励过程，并假设回报可积，有
\[
\begin{aligned}
\mathbb E[G_{t+1}\mid S_t=s]
&=\sum_{s'}\Pr(S_{t+1}=s'\mid S_t=s)\,
  \mathbb E[G_{t+1}\mid S_t=s,S_{t+1}=s']\\
&=\sum_{s'}\Pr(S_{t+1}=s'\mid S_t=s)\,V(s')\\
&=\mathbb E[V(S_{t+1})\mid S_t=s].
\end{aligned}
\]
这里第二个等号用到了马尔可夫性质：\keyterm{已知下一状态是 \(s'\) 后，后续回报的分布不再依赖此前的状态 \(s\)}，因此其条件期望就是从 \(s'\) 出发的价值 \(V(s')\)。

最后利用期望的线性性质，就得到红色箭头处的等式：
\[
\begin{aligned}
\mathbb E[R_t+\gamma G_{t+1}\mid S_t=s]
&=\mathbb E[R_t\mid S_t=s]
  +\gamma\,\mathbb E[G_{t+1}\mid S_t=s]\\
&=\mathbb E[R_t\mid S_t=s]
  +\gamma\,\mathbb E[V(S_{t+1})\mid S_t=s]\\
&=\mathbb E[R_t+\gamma V(S_{t+1})\mid S_t=s].
\end{aligned}
\]
注意，\(V(s')\) 对固定状态 \(s'\) 是一个确定数值，但 \(S_{t+1}\) 尚未确定，所以 \(V(S_{t+1})\) 仍然是随机变量，外面的期望仍然需要保留。

\section*{补充说明：为什么还需要外层期望}

\keyterm{内层的条件期望确实直接等于 \(V(s')\)，但整个式子还要对所有可能的下一状态 \(s'\) 加权平均。}上述推导中的替换发生在这一部分：
\[
\underbrace{\mathbb E[G_{t+1}\mid S_t=s,\ S_{t+1}=s']}_{\text{已经知道下一状态是 }s'}
=V(s').
\]
也就是说，求和号里面的条件期望变成了 \(V(s')\)，但外面的“乘以概率，再求和”仍然保留。

为什么还要这一步？因为原式 \(\mathbb E[G_{t+1}\mid S_t=s]\) \keyterm{只告诉我们当前状态是 \(s\)，没有告诉我们下一状态是什么}。下一步可能到达不同状态，所以要把各状态的价值按到达概率加权：
\[
\mathbb E[G_{t+1}\mid S_t=s]
=
\underbrace{\sum_{s'}\Pr(S_{t+1}=s'\mid S_t=s)\,V(s')}_{\text{各个可能的下一状态的价值，按概率加权平均}}
=
\mathbb E[V(S_{t+1})\mid S_t=s].
\]
这里 \(V(s')\) 是某个指定状态的确定价值，而 \(V(S_{t+1})\) 是“实际到达的下一状态的价值”；由于下一状态尚未确定，它也尚未确定。

例如，下一步各以一半的概率到达 \(A\) 或 \(B\)，且 \(V(A)=5,\ V(B)=3\)。那么 \(V(S_{t+1})\) 可能是 \(5\)，也可能是 \(3\)；但在当前状态 \(s\) 看来，它的期望是 \(\tfrac12\times5+\tfrac12\times3=4\)。因此，\(\mathbb E[G_{t+1}\mid S_t=s]\) 等于确定的平均值 \(4\)，不能直接等于那个可能取 \(5\) 或 \(3\) 的随机变量 \(V(S_{t+1})\)。

\end{document}
